\clearpage
\VerdierPhysicalPageAnchor{0165}%
\begingroup\small\setlength{\parskip}{0pt}\setlength{\emergencystretch}{1.5em}
\VerdierSemanticAnchor{ch4-4-2-1-inclusions}
On a ainsi défini deux foncteurs de $\mathcal{F}\ell_2(\mathcal J)$ dans
$\mathcal A$. On a par définition des inclusions :
\[
B\!\left[\begin{smallmatrix}n\\f,g\end{smallmatrix}\right]\subset H^n(f),
\qquad
Z\!\left[\begin{smallmatrix}n\\f,g\end{smallmatrix}\right]\subset H^n(g).
\]
Soit
$\omega_3=(\mathord\bullet\xrightarrow f\mathord\bullet\xrightarrow g
\mathord\bullet\xrightarrow h\mathord\bullet)$ un objet de
$\mathcal{F}\ell_3(\mathcal J)$.

\VerdierSemanticAnchor{ch4-4-2-2}
\medskip\noindent\textbf{Proposition 4.2.2.}\quad
\textit{On a les inclusions :}
\[
B\!\left[\begin{smallmatrix}n\\g,h\end{smallmatrix}\right]
\subset
Z\!\left[\begin{smallmatrix}n\\f,g\end{smallmatrix}\right]
\subset H^n(g).
\]
La deuxième inclusion est mise pour mémoire. Pour démontrer la première
inclusion, il suffit de montrer que le composé des deux morphismes de la
suite
\VerdierSemanticAnchor{ch4-4-2-2-proof-sequence}
\[
H^{n-1}(h)\xrightarrow{\delta_1}H^n(g)
\xrightarrow{\delta_2}H^{n+1}(f)
\]
est nul. Or l'objet
$(\mathord\bullet\xrightarrow g\mathord\bullet\xrightarrow h\mathord\bullet)$
de $\mathcal{F}\ell_2(\mathcal J)$ fournit la suite exacte
\[
H^{n-1}(h)\xrightarrow{\delta_1}H^n(g)
\xrightarrow u H^n(hg),
\]
et le morphisme de $\mathcal{F}\ell_2(\mathcal J)$
\VerdierSemanticAnchor{ch4-4-2-2-fl2-morphism}
\[
\begin{tikzcd}[column sep=4em,row sep=2.2em]
{}\arrow[r,"f"]\arrow[d,"\operatorname{id}"']&{}
 \arrow[r,"g"]\arrow[d,"\operatorname{id}"']&{}\arrow[d,"h"]\\
{}\arrow[r,"f"']&{}\arrow[r,"hg"']&{}
\end{tikzcd}
\]
fournit le diagramme commutatif
\VerdierSemanticAnchor{ch4-4-2-2-commutative-triangle}
\[
\begin{tikzcd}[column sep=5em,row sep=2.7em]
H^n(g)\arrow[r,"\delta_2"]\arrow[d,"u"']&H^{n+1}(f)\\
H^n(hg)\arrow[ur,"\delta_3"']&{}
\end{tikzcd}
\]
D'où la proposition.
\endgroup
